\documentclass[../../main.tex]{subfiles}%
\graphicspath{{\subfix{./image/}}{\subfix{../../image/}}}

\begin{document}

\section{参考文献与注记}\label{section:ch01-10}

本章篇幅较长，主要内容遵循 Federer\cite{federer1969}，并从 Simon\cite{simon1984} 第 1 章与 Hardt\cite{hardt1979} 补充了材料。

\cref{section:ch01-01}中的\crefrange{theorem:ch01-1-1}{theorem:ch01-1-6}取自\cite{federer1969} 第 2.1.3、2.1.5 节。% XREF p=93 "Theorems 1.1–1.6 in Section 1.1 are from [F, Sections 2.1.3, 2.1.5]" | p=93 "Theorems 1.1–1.6 in Section 1.1 are from [F, Sections 2.1.3, 2.1.5]"

$\pi$-$\lambda$ 定理（\cref{theorem:ch01-1-4}）来自 Durrett\cite{durrett2005} 与 Stroock\cite{stroock1999}。% XREF p=94 "The π-λ Theorem is from Durrett [D] and Stroock [Sk]"
\cref{lemma:ch01-1-1}取自\cite{federer1969} 第 2.2.2 节，\cref{theorem:ch01-1-8}取自\cite{federer1969} 第 2.2.5 节。% XREF p=94 "Lemma 1.1 is [F, Section 2.2.2]" | p=94 "and Theorem 1.8 is [F, Section 2.2.5]"
\cref{theorem:ch01-1-9}（Caratheodory 判据）与\cref{theorem:ch01-1-11}取自\cite{federer1969} 第 2.3.2 节；\cref{theorem:ch01-1-12}取自\cite{federer1969} 第 2.3.3 节。% XREF p=94 "Theorem 1.9 (Caratheodory's criterion) and Theorem 1.11 are from [F, Section 2.3.2]" | p=94 "Theorem 1.9 (Caratheodory's criterion) and Theorem 1.11 are from [F, Section 2.3.2]" | p=94 "Theorem 1.12 is [F, Section 2.3.3]"

\cref{section:ch01-02}中的\cref{theorem:ch01-1-13}是\cite{federer1969} 第 3.1.13、3.1.14 节的简化版本。% XREF p=94 "Theorem 1.13 in Section 1.2 is a simplified version of [F, Sections 3.1.13, 3.1.14]" | p=94 "Theorem 1.13 in Section 1.2 is a simplified version of [F, Sections 3.1.13, 3.1.14]"
Lusin 定理（\cref{theorem:ch01-1-14}）与 Egoroff 定理（\cref{theorem:ch01-1-16}）见\cite{federer1969} 第 2.3.5、2.3.7 节。% XREF p=94 "Lusin's and Egoroff's theorems are in [F, Sections 2.3.5, 2.3.7]" | p=94 "Lusin's and Egoroff's theorems are in [F, Sections 2.3.5, 2.3.7]"
一般的积分理论可在\cite{federer1969} 第 2.4 节中找到；Fatou 引理（\cref{theorem:ch01-1-17}）、单调收敛定理（\cref{theorem:ch01-1-18}）与控制收敛定理（\cref{theorem:ch01-1-19}）特别见\cite{federer1969} 第 2.4.6、2.4.7、2.4.9 节。% XREF p=94 "The general theory of integration may be found in [F, Section 2.4]" | p=94 "see in particular [F, Sections 2.4.6, 2.4.7, 2.4.9] for Fatou's Lemma" | p=94 "the Monotone Convergence Theorem" | p=94 "and the Dominated Convergence Theorem"
我们对乘积测度与 Fubini 定理（\cref{theorem:ch01-1-22}）的处理直接取自\cite{federer1969} 第 2.6 节。% XREF p=94 "Our treatment of product measure and Fubini's Theorem is taken directly from [F, Section 2.6]"
\cref{theorem:ch01-1-4}为 Fubini 定理（\cref{theorem:ch01-1-22}）的断言 (iii) 提供了另一条途径；见 Durrett\cite{durrett2005} 或 Stroock\cite{stroock1999}。% XREF p=94 "The π-λ Theorem provides an alternative approach to assertion (iii) of Fubini's Theorem" | p=94 "assertion (iii) of Fubini's Theorem; see Durrett [D] or Stroock [Sk]"

我们在 Vitali 覆盖定理（\cref{theorem:ch01-1-24}）上主要依靠 Simon\cite{simon1984}，在 Besicovitch 覆盖定理（\cref{theorem:ch01-1-27}）上主要依靠 Hardt\cite{hardt1979}。% XREF p=94 "We relied heavily on Simon [S] for Vitali's Covering Theorem" | p=94 "and Hardt [H] for Besicovitch's Covering Theorem"
另可参看\cite{federer1969} 第 2.9.12、2.9.13 节。Furedi 与 Loeb\cite{furedi1994} 讨论了 Besicovitch 定理的最佳常数。\cref{section:ch01-06}中的微分理论基于\cite{hardt1979}、\cite{simon1984} 与\cite{federer1969} 第 2.9 节。近似极限与近似连续见\cite{federer1969} 第 2.9.12、2.9.13 节。% XREF p=94 "The differentiation theory in Section 1.6 is based on [H], [S], and [F, Section 2.9]"

Riesz 表示定理（\cref{theorem:ch01-1-38}）的证明取自\cite{simon1984} 第 4.1、4.2 节（参看\cite{federer1969} 第 2.5 节）。% XREF p=94 "We took the proof of the Riesz Representation Theorem from [S, Sections 4.1, 4.2] (cf. [F, Section 2.5])"
A. Damlamian 向我们展示了\cref{section:ch01-08}中\cref{theorem:ch01-1-39}的证明。% XREF p=94 "A. Damlamian showed us the proof of Theorem 1.39 in Section 1.8" | p=94 "A. Damlamian showed us the proof of Theorem 1.39 in Section 1.8"
关于\cref{section:ch01-09}中的\cref{theorem:ch01-1-41}，见 Giusti\cite{giusti1984} 附录 A。% XREF p=94 "See Giusti [G, Appendix A] for Theorem 1.41 in Section 1.9" | p=94 "See Giusti [G, Appendix A] for Theorem 1.41 in Section 1.9"
\cref{theorem:ch01-1-43}有时称为 Dunford--Pettis 定理。% XREF p=94 "Theorem 1.43 is sometimes called the Dunford-Pettis Theorem"
咬去引理（\cref{theorem:ch01-1-44}）出现在 Brooks--Chacon\cite{brooks1980} 中，我们的证明遵循 Ball--Murat\cite{ball1989}。% XREF p=94 "The Biting Lemma appears in Brooks-Chacon [B-C] and our proof follows Ball-Murat [B-M]"
L. Tartar 把 Young 测度方法引入偏微分方程理论；例如见\cite{tartar1979}。\cref{theorem:ch01-1-45}与\cref{theorem:ch01-1-46}的证明取自\cite{evans1990}。% XREF p=94 "The proofs for Theorems 1.45 and 1.46 are from [E1]" | p=94 "The proofs for Theorems 1.45 and 1.46 are from [E1]"

测度论的优秀入门书很多。我们特别推荐 DiBenedetto\cite{dibenedetto2002}、Fitzpatrick--Royden\cite{fitzpatrick2010}、Folland\cite{folland1984} 与 Wheeden--Zygmund\cite{wheeden1977}。Krantz--Parks\cite{krantz1999} 与 Lin--Yang\cite{lin2002} 是两本很好的更高阶教材。Oxtoby\cite{oxtoby1980} 对测度论与 Baire 范畴在刻画「可忽略」集合上的对照作了引人入胜的讨论。

\end{document}
